\documentclass[11pt]{article}
\usepackage{coursenotes}
\begin{document}
\notesheader{4}{Causal Trees and Forests}{Searching for heterogeneity without fooling yourself}

\begin{decision}
A search through customer fields finds a segment that responds three times as much as average. Should the
firm build a campaign on it? Only if the segment's effect is measured on data that did not choose it. These
notes show why the search inflates what it finds, how honesty removes the inflation, and how forests turn the
same idea into an estimate for every customer.
\end{decision}

\begin{goals}
  \item explain why a causal splitting criterion ignores prognostic splits while an outcome tree favors them;
  \item quantify the winner's curse and show that honest estimation removes it;
  \item run the discover--estimate--recommend protocol;
  \item describe a causal forest as adaptive neighbor weights with local centering, and state what its
  intervals require;
  \item test whether a CATE model has found real heterogeneity (calibration test and GATES).
\end{goals}

%=====================================================================================
\sectionone

\subsection{An estimand chosen by the data}

Units have covariates $X$, a randomized treatment $D$ with known $\Pr(D=1) = e$, and outcome $Y$. A
\emph{partition} $\Pi = \{\ell_1, \dots, \ell_L\}$ divides covariate space into \emph{leaves}, and the
\emph{leaf CATE} is $\tau_\ell = \E[Y(1) - Y(0) \mid X \in \ell]$. When the leaves are the output of a search, the
estimand itself depends on the data that found it. Estimand, identification and estimation still apply, but
only in order: the partition must be fixed before the data used to estimate it are consulted. The leaf
estimate on a sample $S$ is $\hatt_\ell(S) = \bar Y_{1,\ell}(S) - \bar Y_{0,\ell}(S)$.

\subsection{Splitting on the effect}

A regression tree (CART) splits a parent $P$ into children $L, R$ to maximize
$\tfrac{n_L}{n_P}(\bar Y_L - \bar Y_P)^2 + \tfrac{n_R}{n_P}(\bar Y_R - \bar Y_P)^2$. A \emph{causal tree} replaces outcome
means with effect estimates:
\[
  \Delta = \tfrac{n_L}{n_P}(\hatt_L - \hatt_P)^2 + \tfrac{n_R}{n_P}(\hatt_R - \hatt_P)^2
  = \frac{n_L n_R}{n_P^2}(\hatt_L - \hatt_R)^2 ,
\]
the second form holding when $\hatt_P$ is the size-weighted average of the children. Each leaf needs both
treated and control units.

\begin{prop}{A causal criterion ignores prognostic splits}{prognostic}
If a split leaves the CATE unchanged ($\tau_L = \tau_R$), its population $\Delta$ is zero whatever it does to
outcomes. Under randomization with $e = \tfrac12$, the outcome gap between the two sides of a split on $V$ is
\[
  \E[Y \mid V=1] - \E[Y \mid V=0] = \{\mu_0(1) - \mu_0(0)\} + \tfrac12\{\tau(1) - \tau(0)\}.
\]
\end{prop}
\begin{proof}
With $\tau_L = \tau_R$ both squared terms vanish. For the outcome gap,
$\E[Y \mid V=v] = \tfrac12\E[Y(1) \mid v] + \tfrac12\E[Y(0) \mid v] = \mu_0(v) + \tfrac12\tau(v)$; subtract.
\end{proof}

Effect heterogeneity enters an outcome split at half weight, on top of whatever baseline difference the variable
carries. Outcome trees therefore spend their depth on prognostic variables (spend, tenure) before they reach
effect modifiers.

\subsection{The winner's curse}

\begin{prop}{Selection inflates the winner}{curse}
Let $\hat T_1, \dots, \hat T_k$ be unbiased for $\tau_1, \dots, \tau_k$ and $j^\ast = \arg\max_j \hat T_j$. Then
$\E[\hat T_{j^\ast} - \tau_{j^\ast}] \ge 0$. If the estimates are independent $N(\tau, s^2)$ with a common $\tau$,
$\E[\hat T_{j^\ast}] = \tau + s\,\E[\max(Z_1, \dots, Z_k)]$ with $Z_j$ standard normal.
\end{prop}
\begin{proof}
$\max_j \hat T_j \ge \hat T_m$ for every $m$, so $\E[\hat T_{j^\ast}] \ge \max_m \tau_m \ge \E[\tau_{j^\ast}]$. With a common
$\tau$, $\hat T_j = \tau + sZ_j$.
\end{proof}

$\E[\max Z_j]$ is $1.03$ for $k = 4$, $1.54$ for $k = 10$, $1.87$ for $k = 20$ and $2.25$ for $k = 50$; for large
$k$ it grows like $\sqrt{2\ln k}$. A tree search over fields and cut points compares hundreds of candidates. The winner's
effect, reported on the data that chose it, is too high, and its standard error is too small because it ignores
the search.

\subsection{Honesty}

\begin{prop}{Honest leaf estimates are unbiased}{honest}
Split the data at random into $A$ and $B$, let the partition be any function of $A$ alone, and let $B$ contain its
own randomized assignment. Then for every leaf with both arms present in $B$,
$\E[\hatt_\ell(B) \mid \Pi] = \tau_\ell$, and the difference-in-means standard error computed on $B$ is valid given
$\Pi$.
\end{prop}
\begin{proof}
$B$ is independent of $A$ and hence of $\Pi$, so conditioning on $\Pi$ leaves $B$ unchanged: within a leaf, $B$'s
units are a random sample from that leaf with randomized treatment, and the experiment results of Meetings 1
and 2 apply.
\end{proof}

The proof uses one fact: the partition does not depend on $B$. Consulting $B$ while tuning the tree, even
informally, turns it into discovery data. The cost is variance, because each half does only one job. If a leaf's full-sample
SE is $s$, honesty raises its variance from $s^2$ to $2s^2$; reporting the winner from the search instead adds a
squared bias of about $(s\,\E[\max Z_j])^2$, which exceeds $s^2$ once four or more equal candidates are compared.
Splitting also spends sample size: a leaf with 100 per arm and rates $0.30$ vs $0.10$ has SE $0.055$; each half of it
has SE $0.077$.

\paragraph{The protocol.} \textbf{H1 Discover} on $A$ by any method and \emph{commit}: write the leaves and
the intended decision per leaf. \textbf{H2 Estimate} each committed leaf's effect and SE on $B$, with what the
decision needs (the treated-arm rate, for a per-redemption break-even). \textbf{H3 Recommend} by placing each
interval against the leaf's own threshold: above (act), straddling (undecided), below (do not act). With $L$
leaves, a Bonferroni level $1 - 0.05/L$ controls the chance that any line on the card is wrong.

\subsection{From trees to forests}

A single tree is unstable. A \emph{causal forest} grows many honest trees, each on a random subsample drawn
without replacement with a random subset of candidate variables at each split. It is best read as a way of
choosing, for each target $x$, which training units are its neighbors.

\begin{prop}{The forest estimate as a weighted residual regression}{forest}
For a target $x$, tree $b$ places $x$ in leaf $L_b(x)$. Define
\[
  \alpha_i(x) = \frac1B\sum_{b=1}^B \frac{\ind{X_i \in L_b(x)}}{|L_b(x)|}, \qquad \textstyle\sum_i \alpha_i(x) = 1 .
\]
The forest solves $\sum_i \alpha_i(x)\big[(Y_i - \hat\ell(X_i)) - (D_i - \hat e(X_i))\,\theta\big](D_i - \hat e(X_i)) = 0$:
\[
  \hatt(x) = \frac{\sum_i \alpha_i(x)\,(Y_i - \hat\ell(X_i))(D_i - \hat e(X_i))}{\sum_i \alpha_i(x)\,(D_i - \hat e(X_i))^2},
\]
with $\hat\ell$ and $\hat e$ estimated out of bag.
\end{prop}
\begin{proof}
The estimating equation is the first-order condition of a weighted least-squares fit of the outcome residual on the
treatment residual, the R-learner's loss (Meeting~3) localized by $\alpha_i(x)$; solve the linear equation for
$\theta$.
\end{proof}

Subtracting $\hat\ell$ and $\hat e$ first is \emph{local centering}: it removes what prognostic variables explain, so
splits concentrate on effect modifiers, and it protects the estimate when treatment probabilities vary with $x$.
The neighbors are adaptive: the trees decide which directions in $x$ matter for the effect.

\paragraph{Intervals.} Under conditions (a Lipschitz $\tau$, covariates with a density bounded away from zero on
bounded support, overlap, honest trees whose leaves shrink in every direction, and subsamples of size $s$ with
$s \to \infty$, $s/n \to 0$), Wager and Athey (2018) show $\hatt(x)$ is asymptotically normal with a variance the
infinitesimal jackknife estimates. Three cautions: the result is pointwise, not a band over all customers;
intervals are usually wide relative to the spread of $\hatt(X_i)$; and that spread overstates the spread of
true effects because it includes noise. Averaging does not make trees independent: with per-tree variance
$\sigma^2$ and pairwise correlation $\rho$, the forest's variance is $\sigma^2[\rho + (1-\rho)/B]$, so more trees
remove simulation noise but not the uncertainty the trees share.

\paragraph{Variable importance is not a test.} Split counts describe the algorithm. Correlated variables share
splits arbitrarily, high-cardinality variables get more chances, and a field can rank high because it proxies
for the real modifier. Treat a high importance as a hypothesis for fresh data.

\subsection{Has the model found real heterogeneity?}

On held-out data, compute each unit's doubly robust score $\phi_i$ (Meeting~3), which satisfies
$\E[\phi_i \mid X_i] = \tau(X_i)$, and the model's prediction $\hatt(X_i)$ fitted on other data.

\begin{prop}{Best linear predictor (calibration) test}{blp}
In the population regression $\phi = \beta_1 + \beta_2(\hatt(X) - \E[\hatt(X)]) + u$,
\[
  \beta_1 = \E[\tau(X)] = \ATE, \qquad \beta_2 = \frac{\Cov(\tau(X), \hatt(X))}{\Var(\hatt(X))}.
\]
\end{prop}
\begin{proof}
With a centered regressor, $\beta_1 = \E[\phi] = \E[\tau(X)]$ and $\beta_2 = \Cov(\phi, \hatt)/\Var(\hatt)$. Since
$\hatt(X)$ is fixed given $X$ and $\E[\phi \mid X] = \tau(X)$, $\Cov(\phi, \hatt) = \Cov(\tau(X), \hatt(X))$.
\end{proof}

So $\beta_2 = 0$ means the score carries no information about effects; $\beta_2 > 0$ means it detects real
heterogeneity; $\beta_2 = 1$ means its predictions are calibrated; $0 < \beta_2 < 1$ means they are right in
direction but too spread out. The test is about \emph{this} score: a non-rejection does not show heterogeneity
is absent.

\paragraph{Sorted group effects (GATES).} Sort held-out units by $\hatt$ into $K$ groups using cut-points fixed in
advance and estimate each group's effect by difference in means (unbiased by Result~\ref{prop:honest}, since the
grouping was fixed on other data). Real heterogeneity shows as effects rising from bottom to top. These are the
numbers a manager can act on.

\begin{pitfall}[searching and reporting]
Reporting a searched segment's effect from the search data; re-drawing leaves after looking at the estimation
half; reading the spread of predicted effects, or variable importance, as evidence of heterogeneity.
\end{pitfall}

\checkpoint{An analyst searches 20 candidate segments, finds one with lift $0.12$ (SE $0.03$) against an overall
lift of $0.06$, and proposes a campaign. What lift would you expect on fresh data if every segment's true effect
were $0.06$, and what should be done before the campaign?}

%=====================================================================================
\sectiontwo

\paragraph{Setting.} QuickBite tests free delivery on a shopper's next order for 4{,}000 shoppers (50/50), split at random into a discovery half and an estimation half of 2{,}000 each. The
outcome is whether the shopper orders within a week. Overall the offer raises ordering by $0.054$.

\paragraph{Step 1: choosing the first split (discovery half).}
\begin{center}
\begin{tabular}{lccc}
\toprule
Candidate & Left ($n$, $\hatt$) & Right ($n$, $\hatt$) & $\tfrac{n_Ln_R}{n_P}(\hatt_L - \hatt_R)^2$ \\
\midrule
A: basket $\ge$ ¥80 & (1{,}200, 0.03) & (800, 0.09) & $1.73$ \\
B: first-time customer & (600, 0.11) & (1{,}400, 0.03) & $2.69$ \\
\bottomrule
\end{tabular}
\end{center}
(Scores use the unnormalized form $n_Ln_R/n_P \cdot (\hatt_L - \hatt_R)^2$, which ranks splits identically.) The tree
splits on first-time customers. Both are consistent with the overall effect: $(600 \times 0.11 + 1{,}400 \times 0.03)/2{,}000 = 0.054$.

\paragraph{Step 2: the curse in numbers.} If the search ends up comparing 10 leaves, all with true effect $0.05$
and SE $0.02$, the reported winner averages $0.05 + 0.02 \times 1.54 = 0.081$, a 62\% overstatement. Re-estimated on
the estimation half, it averages $0.05$.

\paragraph{Step 3: a forest estimate by hand.} Two honest trees put a target shopper $x$ in leaves $\{1,2,3,4\}$ and
$\{1,2,5,6\}$. Units 1 and 2 are in both, with weight $\tfrac12(\tfrac14 + \tfrac14) = 0.25$; units 3--6 get $0.125$.
With $\hat e = 0.5$:
\begin{center}
\begin{tabular}{ccccccc}
\toprule
Unit & $D$ & $Y$ & $\hat\ell$ & $\alpha$ & $(Y - \hat\ell)(D - 0.5)$ & $\alpha \times$ previous \\
\midrule
1 & 1 & 1 & 0.6 & 0.250 & $0.20$ & $0.0500$ \\
2 & 0 & 0 & 0.5 & 0.250 & $0.25$ & $0.0625$ \\
3 & 1 & 1 & 0.4 & 0.125 & $0.30$ & $0.0375$ \\
4 & 0 & 1 & 0.6 & 0.125 & $-0.20$ & $-0.0250$ \\
5 & 1 & 0 & 0.3 & 0.125 & $-0.15$ & $-0.01875$ \\
6 & 0 & 0 & 0.4 & 0.125 & $0.20$ & $0.0250$ \\
\midrule
Sum & & & & 1.000 & & $0.13125$ \\
\bottomrule
\end{tabular}
\end{center}
The denominator is $\sum_i \alpha_i(D_i - 0.5)^2 = 0.25$, so $\hatt(x) = 0.13125/0.25 = 0.525$ (Result~\ref{prop:forest}).
With six units this only illustrates the arithmetic.

\paragraph{Step 4: is the heterogeneity real?} On a held-out set the calibration regression gives $\hat\beta_1 = 0.052$
(SE $0.010$) and $\hat\beta_2 = 0.64$ (SE $0.21$): heterogeneity detected ($z = 3.05$), predictions somewhat too spread
out but not significantly so ($z = (1 - 0.64)/0.21 = 1.71$). GATES by quintile of predicted effect, each with SE
about $0.02$:

\begin{center}
\begin{tikzpicture}
\begin{axis}[ybar, width=0.62\linewidth, height=4.6cm, bar width=16pt, symbolic x coords={Bottom,2,3,4,Top},
  xtick=data, ylabel={group effect}, ymin=-0.03, ymax=0.15, axis lines=left, enlarge x limits=0.15,
  yticklabel style={/pgf/number format/fixed}, nodes near coords, nodes near coords style={font=\scriptsize, yshift=6pt},
  every node near coord/.append style={/pgf/number format/fixed, /pgf/number format/precision=3}]
\addplot[fill=greentint, draw=sbgreen, error bars/.cd, y dir=both, y explicit]
  coordinates {(Bottom,0.010) +- (0,0.039) (2,0.030) +- (0,0.039) (3,0.045) +- (0,0.039) (4,0.068) +- (0,0.039) (Top,0.107) +- (0,0.039)};
\draw[dashed, warnred] (axis cs:Bottom,0.06) -- (axis cs:Top,0.06);
\end{axis}
\end{tikzpicture}

\small Figure 1. Group effects with 95\% intervals; the dashed line is a break-even of $0.06$.
\end{center}

The group effects average $0.052$, matching $\hat\beta_1$. Against a break-even of $0.06$, only the top quintile's
interval lies wholly above it; the bottom quintile's lies wholly below, and quintiles 2 to 4 straddle it. Individual predictions are not precise enough to go
further.

\begin{exercise}[computation]
The analyst wants to split the first-time-customer leaf (300 per arm, ordering rates $0.31$ and $0.20$) again, into two
equal halves. (a) What is the leaf's SE now, and each half's SE if both halves had the same rates? (b) How far apart
would the halves' true effects have to be for the gap to be detected at 5\% with 80\% power?
\end{exercise}
\answer{(a) Now: $\sqrt{0.31 \times 0.69/300 + 0.20 \times 0.80/300} = 0.0353$. Halves: $\sqrt{0.2139/150 + 0.16/150} = 0.0499$.
(b) The gap's SE is $\sqrt{2} \times 0.0499 = 0.0706$, so the detectable gap is $2.8 \times 0.0706 \approx 0.20$, nearly
twice the leaf's whole effect of $0.11$. The split is not worth making.}

\begin{exercise}[judgment]
The forest's variable-importance chart ranks ``app version'' first. A product manager concludes that the app
version drives the response to free delivery. What two analyses would you run before agreeing, and on which data?
\end{exercise}
\answer{(1) GATEs by app version estimated on held-out data that played no part in growing the forest, with the SE
of the difference between versions. (2) Check what app version proxies for (new users are on the newest version):
compare effects by app version within first-time and repeat customers. Importance counts splits; it is not a test,
and correlated or high-cardinality fields absorb splits arbitrarily.}

%=====================================================================================
\sectionthree

\begin{extension}[Generalized random forests]
Any parameter defined by a local moment condition $\E[\psi(O;\theta) \mid X = x] = 0$ (a quantile, a local IV effect)
can be estimated with forest weights by solving $\sum_i \alpha_i(x)\psi(O_i;\theta) = 0$. A gradient trick turns split
selection into ordinary CART on pseudo-outcomes. The \texttt{grf} package implements causal, quantile, IV and survival
forests, the calibration test and doubly robust GATEs (Athey, Tibshirani and Wager, 2019).
\end{extension}

\begin{extension}[Honest criteria, inference on winners, shrinkage]
Athey and Imbens's criterion subtracts the estimation variance a split will create, discouraging small leaves
directly. When a winner must be reported from the data that chose it, conditional intervals (Andrews, Kitagawa
and McCloskey, 2024) are valid given the selection; empirical Bayes shrinkage pulls noisy segments toward the
mean and tempers the winner automatically.
\end{extension}

\subsection*{Annotated reading}
\begin{readinglist}
\reading{Athey and Imbens (2016), ``Recursive partitioning for heterogeneous causal effects'',
\emph{PNAS}}{Causal trees and honesty, with the splitting criterion}{Sections 1--4}{2}
\reading{Wager and Athey (2018), ``Estimation and inference of heterogeneous treatment effects using random
forests'', \emph{JASA}}{The asymptotic normality result and its conditions}{Sections 1--3}{3}
\reading{Athey, Tibshirani and Wager (2019), ``Generalized random forests'', \emph{Annals of Statistics}}{Forests
as weights for any moment condition}{Sections 1--2, 6}{3}
\reading{Chernozhukov, Demirer, Duflo and Fern\'andez-Val (2025), ``Fisher--Schultz Lecture: Generic machine learning inference on
heterogeneous treatment effects in randomized experiments'', \emph{Econometrica}}{The calibration test and GATES with repeated
splitting}{Sections 1--3}{2}
\reading{Andrews, Kitagawa and McCloskey (2024), ``Inference on winners'', \emph{QJE}}{Valid intervals for the
best-looking segment}{Sections 1--2}{3}
\end{readinglist}

\end{document}
